\section{Trusted Data：高风险模型的数据治理}

前一节依赖训练标签和 verified outcome，本节追问这些数据从哪里来、谁能访问、怎样证明合法与可靠。高风险安全数据既稀缺又敏感，不能沿用普通互联网抓取或“先复制到 data lake 再治理”的方式。数据能力的核心壁垒不是文件数量，而是授权、provenance、segregation、label protocol 与持续审计。

\subsection{Provenance、segregation 与 retention}

\term{provenance}（来源与处理谱系）记录数据由谁、在何种授权下、经过哪些变换进入哪个训练或评估版本。\term{data segregation}（数据隔离）把极高敏感数据与一般日志、开发环境、低权限账户和其他租户分开。\term{retention}（保留策略）规定数据为何保留、保存多久、何时删除或进入 legal hold，而不是默认永久保存。

\lecturefigure{06-trusted-data.png}{高风险训练数据需要可信来源、可追溯处理、隔离访问和明确保留策略。}{本地字幕 16:20--20:30；Thorn/NCMEC 合作语境；概念重绘。}

\begin{knowledgebox}{读图：数据治理贯穿整个 lifecycle}
Trusted source 证明授权与用途；provenance 连接原始证据、去标识变换、label 和 dataset version；segregation 限制环境、网络与人员；retention/audit 则处理删除、事件响应和监管检查。任何一块缺失，团队都可能无法解释某个模型为何学到某种模式、某位员工为何能访问数据，或删除请求是否真正覆盖 derived artifacts。
\end{knowledgebox}

\begin{importantbox}{Label quality 不等于多数投票}
高风险标签需要法律定义、平台 policy 与操作指南的明确边界。训练集应区分 verified positive、policy violation、uncertain、benign hard negative 和 unavailable context；评估还要记录 reviewer agreement 与 adjudication。把所有“不舒服”内容混成一个正类，会产生不可解释模型，也会扩大对合法表达的误伤。
\end{importantbox}

\teachervoice{讲者强调，Thorn 能训练这类模型并不是因为找到一个更大的公开数据集，而是通过长期可信合作获得受控、合法且带上下文的材料。课堂提示我们：在敏感领域，data access 本身就是治理责任，而非普通采购优势。}

\subsection{训练、评估与生产数据必须分工}

本节把数据用途拆开，避免模型团队在同一 reference set 上不断调参后仍把结果称为独立测试。训练数据用于拟合参数，validation 数据用于选择阈值和模型，test 数据用于最终评估，production feedback 则需要经过偏差与选择机制分析后才能进入下一版本。对每个集合都应固定 snapshot、版本、访问角色和 contamination 检查。

\begin{warningbox}{Verified feedback 也可能有 selection bias}
只有被模型送去审核的对象更容易获得标签，因此反馈数据天然偏向当前模型认为“可疑”的区域。若直接用这些样本重训，系统可能越来越看不见旧阈值之外的风险。团队需要随机抽样、shadow evaluation、独立 reference set 与外部 red-team evidence 来发现 blind spot。
\end{warningbox}

\subsection{本章小结}

Trusted data 是法律授权、技术隔离、标签协议和审计证据的组合。高风险模型若无法解释训练数据来源、用途和删除路径，即使离线指标很高，也不具备安全部署条件。

